\section{Outstanding challenges}\label{sec:open}
\paragraph{How much of the source is needed for the sign?}
The finite degree formula suggests a first concrete experiment in information reduction. Keep the same circle, time interval and target \(D_F-D_{24}\), but form \(q_0=(G_{26}-G_{24})/G_{24}\) from finite readouts. To predict the sign one must independently bound the rest, \(e\) in \eqref{eq:remaining}, and bound the complete quotient by \(\eta<1\) on the collar. Here \(r(2r-b)=153\). If \(J_0=-\langle(z-c_{24})q_0\rangle\) and the nonlinear allowance stay near the observed values, an error appreciably below \(4.96\times10^{-8}/153\), about \(3.2\times10^{-10}\) in relative units, would suffice. The actual margin must be computed for \(q_0\). Using the complete numerator of Section~\ref{sec:evidence} to determine that margin would spend the information the finite step is intended to save. A meaningful comparison gives a direct evaluator the same finite readouts, real bounds, and residual allowance.

\paragraph{Can the degree estimates be iterated?}
A single signed correction supplies neither a common count at arbitrarily large \(N\) nor a rule for its sign at the next degree. One possible sufficient mechanism is to control \(u_N=\max(0,-D_N)\) on a fixed parity by
\[
 u_{N+2}\le (1-\alpha_N)u_N+\varepsilon_N,\qquad
 0<\alpha_N\le1,\qquad \sum_N\alpha_N=\infty,\qquad
 \varepsilon_N/\alpha_N\longrightarrow0 .
\]
The sums and limits are on that parity and \(\varepsilon_N\ge0\). Indeed, for any \(\delta>0\), eventually \(\varepsilon_N\le\delta\alpha_N\), so \((u_{N+2}-\delta)_+\le(1-\alpha_N)(u_N-\delta)_+\). The product tends to zero and then \(\delta\downarrow0\) gives \(u_N\to0\). To infer \(D_F\ge0\), one must also identify the same limiting cluster and show \(D_N\to D_F\) on its counted chart. A source-dependent inequality of this kind is open. Theorem~\ref{p:oscillation} rules out a general monotonicity shortcut for compact positive sources. A strictly positive lower bound would demand more than vanishing negative part.

\paragraph{Can time transport give new coverage?}
The full forcing in~\eqref{d:weighted-gap} has a signed endpoint representation on every subinterval. Useful backward continuation requires a common two-zero chart, an initial absolute margin, the true exterior weight, and a bound on the worst cumulative loss before a possible event, as in \cite[(4.3)]{agc}. The source must control these jointly if cancellations are to survive. A change of chart needs a verified identification of the enclosed pair, and a cluster of larger multiplicity needs more coordinates. The uniform finite-interval calculation here quantifies approximation error in a region with a known prior count. Extending this information to other heights or longer times, or using it to improve a bound for the de Bruijn--Newman constant, remains open.
